\documentclass[leqno]{article}
% Shared QBT/Soln pack preamble for resources.danieljsmith.org.
\usepackage{graphicx}
\usepackage{dsfont}
\usepackage{mathtools}
\usepackage{amssymb}
\usepackage{amsmath}
\usepackage{mleftright}
\usepackage{amsthm}
\usepackage{float}
\usepackage{ragged2e}
\usepackage{tensor}
\usepackage{fancyhdr}
\usepackage{empheq}
\usepackage[most]{tcolorbox}
\usepackage{enumitem}
\usepackage{dirtytalk}
\usepackage{marginnote}
\usepackage{microtype}
\usepackage[
  a4paper,
  left=0.8in,
  right=1in,
  top=1in,
  bottom=0.5in,
  headheight=14pt,
  headsep=0.4in
]{geometry}
\setlength{\parindent}{0cm}

% TikZ / PGFPlots
\usepackage{pgfplots}
\pgfplotsset{compat=1.18}
\usepgfplotslibrary{fillbetween, polar}
\usetikzlibrary{calc, positioning}
\usetikzlibrary{angles, quotes}
\usetikzlibrary{arrows.meta}
\usetikzlibrary{decorations.pathreplacing, decorations.markings}
\usetikzlibrary{patterns}
\usetikzlibrary{intersections}

% Links and contents. The contents follow the TMUA papers: compact entries,
% no dot leaders, and each question's mark total in grey after its label.
\usepackage{tocloft}
\tocloftpagestyle{fancy}
\setlength{\cftbeforesubsecskip}{2pt}
\setlength{\cftsubsecindent}{1.5em}
\setlength{\cftsubsecnumwidth}{0pt}
\renewcommand{\cftsubsecfont}{\small\color{black}}
\renewcommand{\cftsubsecpagefont}{\small\color{black}}
\renewcommand{\cftsubsecleader}{\hfill}
\usepackage[colorlinks=true,linkcolor=black,citecolor=magenta,urlcolor=black]{hyperref}%
\urlstyle{same}

% Pack macros
% \marks also adds its value to the current question's total in the .aux
% file, so the contents show each total from the second pass. The write
% happens at shipout, so a box typeset more than once counts once. The
% pack's mark total for the contents-page title is summed the same way.
\makeatletter
\newcount\djs@totalmarks
\newcommand{\djsAddMarks}[2]{%
  \global\advance\djs@totalmarks#2\relax
  \@ifundefined{djs@marks@#1}%
    {\expandafter\xdef\csname djs@marks@#1\endcsname{#2}}%
    {\expandafter\xdef\csname djs@marks@#1\endcsname
      {\the\numexpr\csname djs@marks@#1\endcsname+#2\relax}}}
\RenewDocumentCommand{\marks}{o m}{%
  \protected@write\@auxout{}{\string\djsAddMarks{\number\value{enumi}}{#2}}%
  \IfValueTF{#1}%
    {\marginnote{\raggedright\textbf{[#2]}}[#1]}%
    {\marginnote{\raggedright\textbf{[#2]}}}%
}
\makeatother

\newcommand{\vspacerule}[2]{%
  \par\vspace*{#1}%
  \noindent\hrulefill\par
  \vspace*{#2}%
}

\definecolor{scarlet}{RGB}{205, 40, 75}
% Site colours for the contents-page title: the accent, the rule colour, a
% stronger rule colour for the path separator, and the muted text colour.
\definecolor{djsBurgundy}{RGB}{118, 59, 67}
\definecolor{djsHairline}{RGB}{216, 213, 205}
\definecolor{djsSlash}{RGB}{170, 166, 158}
\definecolor{djsMuted}{RGB}{96, 101, 106}

% Worked-solution box, adapted from the TMUA worked-solutions box. A white breakable box with a left scarlet rule and no surrounding
% frame or tint. A part that breaks fills the rest of its page, so the rule
% runs to the bottom margin; the unbroken or last part ends in a short
% horizontal foot rule so a solution visibly stops. Unlike the TMUA box there
% is no answer label. The box does not grow to the left, so the rule sits at
% the text edge; growing to the right keeps the usable line width close to the
% old tcolorbox Solution box.
% Spacing inside the box: paragraphs are 6pt apart, so source blank lines show
% (solutions put one after each display and none between a lead-in and its
% display). Every display gets 6pt above and 10pt below, whatever the length
% of its lead-in line; align rows are spread by an 8pt \jot; solution parts
% are 14pt apart and sub-parts 8pt; \left and \right space like ordinary
% brackets (mleftright); no single line is stranded at a page break; and
% neither the Solution label nor the line before a list ends a page. No
% skip can shrink, so full pages are not squeezed.
\newtcolorbox{djsSolution}{%
  enhanced, breakable, frame hidden, sharp corners, colback=white,
  height fixed for=first and middle, height fill,
  extras unbroken and last={natural height},
  borderline west={2.2pt}{0pt}{scarlet},
  left=11pt, right=0pt, top=5pt, bottom=13pt,
  before skip=13pt, after skip=4pt,
  grow to right by=10mm,
  overlay unbroken and last={%
    \fill[scarlet] (frame.south west)
      rectangle ([xshift=3.5cm, yshift=2.2pt]frame.south west);},
  before upper={{\color{scarlet}\bfseries Solution}\par\nobreak\vspace{4pt}%
    \setlength{\parskip}{6pt plus 2pt}%
    \setlength{\abovedisplayskip}{6pt plus 2pt}%
    \setlength{\abovedisplayshortskip}{6pt plus 2pt}%
    \setlength{\belowdisplayskip}{10pt plus 2pt}%
    \setlength{\belowdisplayshortskip}{10pt plus 2pt}%
    \setlength{\jot}{8pt}%
    \setlist{parsep=6pt plus 2pt, topsep=6pt plus 2pt, beginpenalty=10000}%
    \setlist[enumerate,1]{itemsep=14pt plus 2pt}%
    \setlist[enumerate,2]{itemsep=8pt plus 2pt}%
    \mleftright
    \clubpenalty=10000 \widowpenalty=10000 \displaywidowpenalty=10000\relax}}

% Optional smaller figure inside a solution box. \scalebox shrinks node text
% with the drawing. Not applied to existing figures.
\newcommand{\SolnFigureScale}{0.85}
\newsavebox{\djsSolnFigureBox}
\newenvironment{djsSolutionFigure}
  {\par\vspace{4pt}\begin{center}\begin{lrbox}{\djsSolnFigureBox}}
  {\end{lrbox}\scalebox{\SolnFigureScale}{\usebox{\djsSolnFigureBox}}%
   \end{center}\vspace{2pt}\par}

\newcommand{\cxl}[2]{%
  \tikz[baseline=(X.base)]{%
    \node[inner sep=0pt, outer sep=0pt, text=#1] (X) {$#2$};
    \draw[#1, line width=0.4pt] (X.south west)--(X.north east);
  }%
}

% Page-style hook run at the contents page and each question's opening page.
% A no-op for QBT packs; \djsSolnHeader points it at the fancy header.
\newcommand{\djsQuestionPageStyle}{}

% A contents entry: "Question N" in a box wide enough for two digits, so the
% mark totals line up, then the total in grey. Bookmarks keep "Question N".
\makeatletter
\newlength{\djsTocLabelWidth}
\newcommand{\djsTocMarks}[1]{%
  \@ifundefined{djs@marks@#1}{}{%
    {\footnotesize\itshape\color{black!45}%
     \csname djs@marks@#1\endcsname~mark%
     \ifnum\csname djs@marks@#1\endcsname=1 \else s\fi}}}
\newcommand{\djsTocEntry}[1]{%
  \texorpdfstring
    {\settowidth{\djsTocLabelWidth}{Question 00}%
     \makebox[\dimexpr\djsTocLabelWidth+0.9em\relax][l]{Question #1}%
     \djsTocMarks{#1}}%
    {Question #1}}
\makeatother

% The question count for the contents-page title is written to the .aux at
% the end of the document, so it also appears from the second pass.
\newcounter{djsQuestions}
\newcommand{\questionitem}{%
  \djsQuestionPageStyle
  \item\phantomsection
  \stepcounter{djsQuestions}%
  \addcontentsline{toc}{subsection}{\protect\djsTocEntry{\number\value{enumi}}}%
}
\makeatletter
\newcommand{\djsSetQuestionCount}[1]{\gdef\djs@questioncount{#1}}
\AtEndDocument{\immediate\write\@auxout
  {\string\djsSetQuestionCount{\number\value{djsQuestions}}}}
\makeatother

\renewcommand{\contentsname}{Questions}

% \djsPackHeader[<link content tag>]{<left header>}; the tag says which
% document a click on the header link came from.
\newcommand{\djsPackHeader}[2][questions]{%
  \pagestyle{fancy}%
  \fancyhead{}%
  \fancyfoot{}%
  \renewcommand{\headrulewidth}{0.9pt}%
  \lhead{#2}%
  \chead{}%
  \rhead{\href{https://danieljsmith.org/?utm_source=pdf&utm_campaign=further-maths&utm_content=#1}{danieljsmith.org}}%
}

% The header's topic is also the contents-page title, under which a label
% says which booklet this is.
\newcommand{\djsPackTitle}{}
\newcommand{\djsBookletLabel}{}
\newcommand{\djsQbtHeader}[1]{%
  \djsPackHeader{#1}%
  \renewcommand{\djsPackTitle}{#1}%
  \renewcommand{\djsBookletLabel}{Questions by topic}%
}

% The colophon at the foot of page one, independent of the page style: the
% contents-page title's hairline led by a short scarlet rule, then the date on
% the left and a link to the site on the right, all inside the bottom margin.
% Its sizes are fixed, so it matches the TMUA papers' colophon. The document
% supplies the date; recompiling never changes it.
\makeatletter
\newcommand{\djsLastUpdated}[1]{%
  \AtBeginDocument{\AddToHookNext{shipout/foreground}{\djs@colophon{#1}}}}
\newcommand{\djs@colophon}[1]{%
  \put(0.8in,-\dimexpr\paperheight-0.47in\relax){%
    \makebox[0pt][l]{\color{djsHairline}%
      \rule[0.9pt]{\dimexpr\paperwidth-1.8in\relax}{0.6pt}}%
    \makebox[0pt][l]{\color{scarlet}\rule{3.5cm}{2.2pt}}}%
  \put(0.8in,-\dimexpr\paperheight-0.32in\relax){\makebox[0pt][l]{%
    \fontsize{9}{11}\selectfont Last updated #1}}%
  \put(\dimexpr\paperwidth-1in\relax,-\dimexpr\paperheight-0.32in\relax){%
    \makebox[0pt][r]{\fontsize{7}{8.5}\selectfont\color{djsMuted}%
      Latest version at \href{https://resources.danieljsmith.org/}%
        {\color{black}resources.danieljsmith.org}}}}
\makeatother

% Solutions packs show the header only on the contents page and each
% question's opening page, so it never appears part-way through a solution.
% \thispagestyle marks the next page shipped out, which is the question's
% page because every solution question starts after a \newpage.
\newcommand{\djsSolnHeader}[1]{%
  \djsPackHeader[solutions]{\textit{Solutions} - #1}%
  \pagestyle{empty}%
  \renewcommand{\djsQuestionPageStyle}{\thispagestyle{fancy}}%
  \renewcommand{\djsPackTitle}{#1}%
  \renewcommand{\djsBookletLabel}{Worked solutions}%
}

% \djsSitePath{<strand>}{<topic>}, called once after the header, names where
% the pack sits on the site; the contents-page title prints it above the pack
% title, e.g. CORE PURE / FURTHER CALCULUS. Without it the line is left out.
\makeatletter
\newif\ifdjs@sitepath
\newcommand{\djs@strand}{}
\newcommand{\djs@section}{}
\newcommand{\djsSitePath}[2]{%
  \djs@sitepathtrue
  \renewcommand{\djs@strand}{#1}%
  \renewcommand{\djs@section}{#2}%
}
\makeatother

% Further Mechanics packs call \djsMechanicsPack in their preamble, so that
% \djsFrontMatter prints the default value of g above the contents. The box
% keeps its own space above it (before skip is the whole gap below the title,
% not an addition to it). Without the box, \djsGravityNote is just extra space
% before the contents heading, bringing the gap above it close to the gap
% between the box and the heading on Further Mechanics pages.
\newcommand{\djsGravityNote}{\vspace{14pt}}
\newcommand{\djsMechanicsPack}{%
  \renewcommand{\djsGravityNote}{%
    \begin{tcolorbox}[sharp corners, colback=white, colframe=black,
      boxrule=1pt, left=8pt, right=8pt, top=6pt, bottom=6pt,
      before skip=20pt]
    {\bfseries Instructions}\par\vspace{4pt}
    Unless a question states otherwise, take $g = 9.8\,\mathrm{m\,s^{-2}}$
    wherever a numerical value of $g$ is required. Solutions also give final
    answers for $g = 10\,\mathrm{m\,s^{-2}}$ where these differ.
    \end{tcolorbox}%
    \vspace{16pt}}%
}

% Contents-page title: the site path in spaced burgundy capitals, the pack
% title, a hairline led by a short scarlet rule (the solution box's foot rule),
% then the booklet label with the question and mark counts. The counts come
% from the .aux, so they appear from the second pass.
\makeatletter
\newsavebox{\djs@eyebrowbox}
\newcommand{\djs@caps}[1]{\textls[180]{\MakeUppercase{#1}}}
\newcommand{\djs@eyebrow}{%
  \color{djsBurgundy}\footnotesize\bfseries
  \djs@caps{\djs@strand}%
  {\color{djsSlash}\hspace{0.55em}/\hspace{0.55em}}%
  \djs@caps{\djs@section}}
\newcommand{\djs@dot}{\hspace{0.45em}\textperiodcentered\hspace{0.45em}}
\newcommand{\djs@counts}{%
  \@ifundefined{djs@questioncount}{}{%
    \djs@questioncount~question\ifnum\djs@questioncount=1 \else s\fi
    \ifnum\djs@totalmarks>0
      \djs@dot\the\djs@totalmarks~mark\ifnum\djs@totalmarks=1 \else s\fi
    \fi}}
\newcommand{\djsContentsTitle}{%
  \ifdjs@sitepath
    \sbox\djs@eyebrowbox{\djs@eyebrow}%
    \ifdim\wd\djs@eyebrowbox>\linewidth
      \@latex@warning@no@line{The \string\djsSitePath\space line is wider
        than the text, so it wraps. Shorten its topic}%
    \fi
    {\raggedright\djs@eyebrow\par}%
    \vspace{7pt}%
  \fi
  {\raggedright\fontsize{24.88}{28}\selectfont\djsPackTitle\par}%
  \vspace{9pt}%
  \noindent\makebox[0pt][l]{\color{djsHairline}\rule[0.9pt]{\linewidth}{0.6pt}}%
  {\color{scarlet}\rule{3.5cm}{2.2pt}}\par
  \vspace{5pt}%
  {\small\color{djsMuted}\djsBookletLabel\hfill\djs@counts\par}%
  \vspace{8pt}%
}
\makeatother

\newcommand{\djsFrontMatter}{%
  \djsQuestionPageStyle
  \vspace*{6pt}%
  \djsContentsTitle
  \djsGravityNote
  \tableofcontents
  \newpage
}


\djsQbtHeader{Improper Integrals and Mean Values}
\djsSitePath{Core Pure}{Further Calculus}
\djsLastUpdated{26 October 2026}

\begin{document}
\djsFrontMatter
\begin{enumerate}[
  label=\textbf{\arabic*.},
  leftmargin=*,
  topsep=0pt,
  partopsep=0pt,
  parsep=0pt
]

% ---- Question 1 ----
\questionitem
% [Source: _QBT__Improper_Integrals_and_Mean_Values, Q1]

The function $\mathrm{f}$ is defined by
\[\mathrm{f}(x) = \dfrac{x}{x^2+1}, \qquad 1 \leqslant x \leqslant 3\] \\
Find the mean value of $\mathrm{f}$ over the interval $1 \leqslant x \leqslant 3$. \\

Give your answer in the form $p\ln q$, where $p$ is a rational number and $q$ is an integer. \marks{4} \\

\newpage

% ---- Question 2 ----
\questionitem
% [Source: _QBT__Improper_Integrals_and_Mean_Values, Q2]

Show that
\[\int_{2}^{\infty} \dfrac{x}{(x^{2}+5)^{\frac{3}{2}}}\,\mathrm{d}x = \dfrac{1}{3}\]
You must show the limiting process clearly. \marks{5} \\

\newpage

% ---- Question 3 ----
\questionitem
% [Source: _QBT__Improper_Integrals_and_Mean_Values, Q3]

By first expressing
\[\dfrac{4}{x(x+1)(x+2)}\]
in partial fractions, show that the improper integral
\[\int_1^{\infty} \dfrac{4}{x(x+1)(x+2)}\,\mathrm{d}x\]
converges, and find its exact value in the form $\ln k$, where $k$ is a rational number to be determined. \marks{8} \\

\newpage

% ---- Question 4 ----
\questionitem
% [Source: _QBT__Improper_Integrals_and_Mean_Values, Q4]

Let
\[y = \dfrac{1}{1 - \sin x}\]
By first writing $y$ in terms of $\sec x$ and $\tan x$, find the exact mean value of $y$ with respect to $x$ over the interval $0 \leqslant x \leqslant \frac{1}{6}\pi$. \marks{6} \\

\newpage

% ---- Question 5 ----
\questionitem
% [Source: _QBT__Improper_Integrals_and_Mean_Values, Q5]

The power drawn, $p$ watts, by an electric motor at time $t$ hours after it is switched on can be modelled by the equation
\[p = 40\tanh^2\left(\tfrac{1}{2}t\right), \qquad t \geqslant 0\] \\
Using an appropriate hyperbolic identity, find, according to the model, the \textbf{mean} power drawn by the motor over the first $2$ hours after it is switched on. Give your answer as an exact value in the form $\frac{a}{e^2+b}$, where $a$ and $b$ are integers. \marks{6} \\

\newpage

% ---- Question 6 ----
\questionitem
% [Source: _QBT__Improper_Integrals_and_Mean_Values, Q6]

\begin{enumerate}[label=\textbf{(\alph*)}]
  \item Explain why
  \[\int_{e}^{\infty} \dfrac{\ln x}{x^{3}}\,\mathrm{d}x\]
  is an improper integral. \marks{1} \\
  \item Use integration by parts to find
  \[\int \dfrac{\ln x}{x^{3}}\,\mathrm{d}x\] \marks[-30pt]{3}
  \item Hence show that
  \[\int_{e}^{\infty} \dfrac{\ln x}{x^{3}}\,\mathrm{d}x\]
  converges, and find its exact value. You should show the limiting process used, and state clearly the behaviour of any limit you rely on. \marks{4} \\
\end{enumerate}

\newpage

% ---- Question 7 ----
\questionitem
% [Source: _QBT__Improper_Integrals_and_Mean_Values, Q7]

\begin{enumerate}[label=\textbf{(\alph*)}]
  \item Express $9x^2 - 18x + 25$ in the form $a(x+p)^2 + q$, where $a$, $p$ and $q$ are integers. \marks{2} \\
  \item Hence use a standard integral from the formula book to find
  \[\int \dfrac{1}{\sqrt{9x^2 - 18x + 25}}\,\mathrm{d}x\] \marks[-30pt]{2}
  \item Determine the exact mean value of the function
  \[\mathrm{g}(x) = \dfrac{1}{\sqrt{9x^2 - 18x + 25}}, \qquad 0 \leqslant x \leqslant 2\]
  giving your answer in the form $A \ln B$, where $A$ and $B$ are rational constants with $B$ an integer. \marks{3} \\
\end{enumerate}

\newpage

% ---- Question 8 ----
\questionitem
% [Source: _QBT__Improper_Integrals_and_Mean_Values, Q8]

Let $y = \ln(1 + x)$. \\
\begin{enumerate}[label=\textbf{(\alph*)}]
  \item Show that the mean value of $y$ with respect to $x$ over the interval $0 \leqslant x \leqslant e - 1$ is
  \[\dfrac{1}{e - 1}\] \marks[-30pt]{4}
  \item Find, in exact form, the mean value of $x$ with respect to $y$ over the interval $0 \leqslant y \leqslant 1$. \marks{3} \\
\end{enumerate}

\newpage

% ---- Question 9 ----
\questionitem
% [Source: _QBT__Improper_Integrals_and_Mean_Values, Q9]

In an electrical circuit, the current $I$ amps at time $t$ seconds is modelled by
\[ I = a\sin nt + b\sin 2nt \]
where $a$, $b$ and $n$ are positive constants. The RMS value of the current, in amps, is defined to be the square root of the mean value of $I^2$ over the interval $0 \leqslant t \leqslant \frac{2\pi}{n}$. \\

Show that the RMS value of the current is
\[\sqrt{\dfrac{a^2+b^2}{2}}\]
amps. \marks{7} \\

\newpage

% ---- Question 10 ----
\questionitem
% [Source: _QBT__Improper_Integrals_and_Mean_Values, Q10]

Prove that
\[\int_{1}^{\infty} \dfrac{x+3}{x^{2}(x+1)}\,\mathrm{d}x = 3 - \ln 4\]
Show the limiting process clearly. \marks{8} \\

\newpage

% ---- Question 11 ----
\questionitem
% [Source: _QBT__Improper_Integrals_and_Mean_Values, Q11]

\begin{enumerate}[label=\textbf{(\alph*)}]
  \item Explain why
  \[\int_{0}^{\infty} \dfrac{10}{(x+2)(3x+1)}\,\mathrm{d}x\]
  is an improper integral. \marks{1} \\
  \item Show that
  \[\int_{0}^{\infty} \dfrac{10}{(x+2)(3x+1)}\,\mathrm{d}x = a\ln b\]
  where $a$ is an integer and $b$ is a rational number, both to be determined. You must show clearly how you deal with the infinite limit. \marks{6} \\
\end{enumerate}

\newpage

% ---- Question 12 ----
\questionitem
% [Source: _QBT__Improper_Integrals_and_Mean_Values, Q12]

\begin{enumerate}[label=\textbf{(\alph*)}]
  \item Explain why
  \[\int_0^3 \dfrac{x}{\sqrt[3]{x^2-1}}\,\mathrm{d}x\]
  is an improper integral. \marks{1} \\
  \item Use an appropriate limit argument to show that this integral converges, and find its exact value. \marks{5} \\
\end{enumerate}

\newpage

% ---- Question 13 ----
\questionitem
% [Source: _QBT__Improper_Integrals_and_Mean_Values, Q13]

\begin{enumerate}[label=\textbf{(\alph*)}]
  \item Explain why
  \[\int_0^{e} \dfrac{(\ln x)^2}{\sqrt{x}}\,\mathrm{d}x\]
  is an improper integral. \marks{1} \\
  \item By using integration by parts twice, show that
  \[\int x^{-\frac{1}{2}} (\ln x)^2\,\mathrm{d}x = 2\sqrt{x}\left((\ln x)^2 - 4\ln x + 8\right) + c\] \marks[-30pt]{4}
  \item You may assume that, for any constant $k > 0$ and any positive integer $n$, $x^k(\ln x)^n \to 0$ as $x \to 0^{+}$. \\

  Show that
  \[\int_0^{e} \dfrac{(\ln x)^2}{\sqrt{x}}\,\mathrm{d}x\]
  exists and find its exact value. \marks{4} \\
\end{enumerate}

\newpage

% ---- Question 14 ----
\questionitem
% [Source: _QBT__Improper_Integrals_and_Mean_Values, Q14]

\begin{enumerate}[label=\textbf{(\alph*)}]
  \item Find
  \[\int_1^a \dfrac{(\ln x)^2}{x^2}\,\mathrm{d}x\]
  where $a > 1$, giving your answer in terms of $a$. \marks{6} \\
  \item Write down the value of
  \[\lim_{a \to \infty} \dfrac{(\ln a)^k}{a}\]
  where $k$ is a positive constant. \marks{1} \\
  \item Hence find
  \[\int_1^{\infty} \dfrac{(\ln x)^2}{x^2}\,\mathrm{d}x\] \marks[-30pt]{2}
\end{enumerate}

\newpage

% ---- Question 15 ----
\questionitem
% [Source: _QBT__Improper_Integrals_and_Mean_Values, Q15]

Show that
\[\int_0^{\infty} \dfrac{5x-3}{(x^2+3)(x+1)}\,\mathrm{d}x = p\sqrt{3}\,\pi + \ln q\]
where $p$ and $q$ are rational numbers to be found. \marks{8} \\

\newpage

% ---- Question 16 ----
\questionitem
% [Source: _QBT__Improper_Integrals_and_Mean_Values, Q16]

\begin{enumerate}[label=\textbf{(\alph*)}]
  \item Explain why
  \[\int_{0}^{1} (\ln y)^2\,\mathrm{d}y\]
  is an improper integral. \marks{1} \\
  \item
  \begin{enumerate}[label=(\roman*)]
    \item Show that the substitution $y = e^{-x}$ transforms
    \[\int_{0}^{1} (\ln y)^2\,\mathrm{d}y \quad \text{into} \quad \int_{0}^{\infty} x^2 e^{-x}\,\mathrm{d}x\]
    explaining clearly how the limits of integration are obtained. \marks{3} \\
    \item Evaluate
    \[\int_{0}^{\infty} x^2 e^{-x}\,\mathrm{d}x\]
    showing the limiting process used. \marks{6} \\
    \item Hence write down the value of
    \[\int_{0}^{1} (\ln y)^2\,\mathrm{d}y\] \marks[-30pt]{1}
  \end{enumerate}
\end{enumerate}

\end{enumerate}
\end{document}
